\documentclass[11pt]{article}
\usepackage[a4paper,margin=27mm]{geometry}
\usepackage[T1]{fontenc}
\usepackage{lmodern,microtype,amsmath,amssymb,amsthm,mathtools,enumitem}
\usepackage[hidelinks]{hyperref}
\newtheorem{theorem}{Theorem}
\newtheorem{lemma}[theorem]{Lemma}
\newtheorem{proposition}[theorem]{Proposition}
\newtheorem{corollary}[theorem]{Corollary}
\theoremstyle{definition}\newtheorem{definition}[theorem]{Definition}
\theoremstyle{remark}\newtheorem{remark}[theorem]{Remark}
\newcommand{\R}{\mathbb R}
\newcommand{\Dscr}{\mathcal D}
\newcommand{\supp}{\operatorname{supp}}
\newcommand{\on}{\mathrm{on}}
\newcommand{\off}{\mathrm{off}}
\newcommand{\esssup}{\operatorname*{ess\,sup}}
\newcommand{\essinf}{\operatorname*{ess\,inf}}
\newcommand{\status}[1]{\par\noindent\emph{Status:} #1\par}
\title{Separation-Gated Evaluators in Two Worlds}
\author{}\date{20 October 2026}
\begin{document}\maketitle

\noindent
This note answers the mechanism-design question left open by the five-region
construction (\emph{Disagreement-Selected Evaluators}, 20 October). There,
evaluators chosen from the disagreement produced genuinely non-additive
comparisons with strict (G), (IB) and (TC), but (S) failed: some rivals were
defended at no world. Following audit's recommendation, the construction below
is designed around (S) from the start.

\begin{enumerate}[leftmargin=*,itemsep=3pt]
\item \emph{Construction} (Definition~\ref{def:rule}). Each world compares two
models only when their disagreement is \emph{separated}: in a fixed priority
order of rank steps, every rank step that the second model adds (a loss for that
world) comes no later than every rank step that it removes (a gain). Inside this
gate a record-based verdict decides. Every other mixed disagreement is left
incomparable.
\item \emph{Main result} (Theorem~\ref{thm:main}). For any continuous monotone
family of interval models through the target, and any continuous verdict with
two sign conditions and a covering condition, the system satisfies reflexivity,
strict (G), (IB), (TC) (in fact it has no cycles at all), (LA), (S) with every
selected model undominated, direct (C), and (LF) and (SC) on the full
four-advice domain. Every non-selected model is dominated by a selected one.
\item \emph{A non-additive instance} (Proposition~\ref{prop:instance}). The
author's two-measure rule (Proposition~1 of the design memo), in rank-step form
and in its ``sticky'' direction (an added rank step costs twice what a removed
one gains), violates (TC) on its own but satisfies every axiom once it is gated,
at both worlds. No integrable strictly guiding field, and no finitely additive
weight positive on non-null sets, endorses the comparisons of either world out
of selected models; the failure already occurs near almost every selected
boundary point (a dense set of indices).
\item \emph{What carries the result, and why two worlds are special}
(Section~\ref{sec:three}). The gate gives (TC) for free; (S) survives the gate
because, in two worlds, the selected models are threshold sets of the priority
order. In three worlds a gate of this strong form makes (S) fail at every
interior index whose selected model is a threshold set at no world
(Proposition~\ref{prop:three}); that this happens for every selection with a
two-dimensional image is CONJECTURE. So the two-world laboratory does not
discriminate between verdicts; the three-world question is where the content is.
\end{enumerate}

\section{Setting}\label{sec:setting}

Two worlds. Problems are pairs $\{g,-g\}$, $g\in\R^2\setminus\{0\}$, with the
normalized Gaussian measure $\mu$ on representatives. A problem is
\emph{unanimous} if $g_1,g_2$ have the same sign and \emph{mixed} otherwise; a
mixed problem has a representative $g=r(1-t,-t)$ with $r>0$, $t\in(0,1)$. The
law of $t$ under $\mu$ has a positive density, so level sets of $t$ are null.
A model is a measurable advice map (accept $g$, silence, accept $-g$, accept
both), modulo null sets; the domain is the full four-advice domain. At world
$i$ the advice at a problem has the rank $r_i\in\{0,1,2,3\}$ (good, silent,
bad, contradictory). On a mixed problem accepting $g=r(1-t,-t)$ is good at
world~$1$ and bad at world~$2$.

For $m=(m_1,1-m_1)$, $m\cdot g=r(m_1-t)$. Hence the coherent model with credal
set $\{m:m_1\in[a,b]\}$, written $\Dscr_{[a,b]}$, accepts $g$ on $\{t<a\}$, is
silent on $\{a\le t\le b\}$, accepts $-g$ on $\{t>b\}$, and is correct on
unanimous problems.

\paragraph{The selected family.} Let $l,u\colon[0,1]\to[0,1]$ be continuous and
strictly increasing with $l\le u$, $l(0)=u(0)=0$ and $l(1)=u(1)=1$. Index the
simplex by $p=p_1\in[0,1]$ and put $C_p\coloneqq\Dscr_{[l(p),u(p)]}$. Then
$C_1=D_1$ and $C_0=D_2$, the ideal models. For audit's target
$K^*=\{m_1\in[\frac13,\frac23]\}$ one may take, for instance,
$l(p)=\frac23p$, $u(p)=\frac43p$ on $[0,\frac12]$ and
$l(p)=\frac13+\frac43(p-\frac12)$, $u(p)=\frac23+\frac23(p-\frac12)$ on
$[\frac12,1]$; then $C_{1/2}=\Dscr_{K^*}$. Since $l,u$ are continuous and level
sets of $t$ are null, $\mu(C_p\,\triangle\,C_q)\to0$ as $q\to p$, where
$\triangle$ is taken over the set where the advice differs.

\paragraph{Rank steps and keys.} For world $i$ let
$I\coloneqq\{(g,j):g\text{ a problem},\ j\in\{1,2,3\}\}$ with the measure
$\mu\otimes$counting, and for a model $X$ put
\[
 U_i(X)\coloneqq\{(g,j)\in I:r_i(g,X)\ge j\},
\]
the set of \emph{world-$i$ rank steps} of $X$. A pointwise weak world-$i$
worsening from $X$ to $Y$ is exactly $U_i(X)\subseteq U_i(Y)$. For a pair
$(X,Y)$ put
\[
 T^\on_i(X,Y)\coloneqq U_i(Y)\setminus U_i(X),\qquad
 T^\off_i(X,Y)\coloneqq U_i(X)\setminus U_i(Y):
\]
the steps that $Y$ adds (world-$i$ losses in passing from $X$ to $Y$) and the
steps that $Y$ removes (world-$i$ gains). Both are determined by the advice of
$X$ and $Y$ on the set where they differ.

Fix a measurable $\tau\colon\{\text{problems}\}\to(0,1)$ with null level sets
(for instance a strictly monotone function of the angle of $g$). Define the
\emph{keys}
\[
 k_1(g,1)\coloneqq l^{-1}(t),\quad k_1(g,2)\coloneqq u^{-1}(t),\qquad
 k_2(g,1)\coloneqq 1-u^{-1}(t),\quad k_2(g,2)\coloneqq 1-l^{-1}(t)
\]
for mixed $g$, and $k_i(g,j)\coloneqq-1-\tau(g)$ in all other cases ($j=3$, or
$g$ unanimous). Every level set of $k_i$ is null.

\begin{lemma}\label{lem:threshold}
For every $p\in[0,1]$, up to null sets,
$U_1(C_p)=\{k_1\ge p\}$ and $U_2(C_p)=\{k_2\ge1-p\}$.
\end{lemma}
\begin{proof}
On a mixed problem, $r_1(g,C_p)$ is $0$ on $\{t<l(p)\}$, $1$ on
$\{l(p)\le t\le u(p)\}$ and $2$ on $\{t>u(p)\}$. So $(g,1)\in U_1(C_p)$ iff
$t\ge l(p)$ iff $k_1(g,1)\ge p$, and $(g,2)\in U_1(C_p)$ iff $t>u(p)$ iff
$k_1(g,2)>p$; steps with $j=3$ and steps at unanimous problems never belong to
$U_1(C_p)$, and their keys are negative. World~2 is the same with
$r_2(g,C_p)$ equal to $2,1,0$ on the three bands.
\end{proof}

So, in the priority order given by the key, every selected model is a
\emph{threshold set}: its world-$i$ steps are exactly the steps of high key.

\section{The gated rule}\label{sec:rule}

\begin{definition}\label{def:rule}
A pair $(X,Y)$ is \emph{$i$-separated} if $T^\on_i(X,Y)$ is non-null and
\[
 \esssup_{T^\on_i(X,Y)}k_i\ \le\ \essinf_{T^\off_i(X,Y)}k_i
\]
(vacuous if $T^\off_i$ is null). Let $\Phi_1,\Phi_2$ be real functions of
pairs of models that depend only on the advice of the two models on the set
where they differ. Define
\[
 X\preceq_iY\quad:\Longleftrightarrow\quad X=Y,\ \text{or}\ T^\off_i(X,Y)
 \text{ is null},\ \text{or}\ \bigl[(X,Y)\text{ is $i$-separated and }
 \Phi_i(X,Y)\ge0\bigr].
\]
\end{definition}

The three clauses are: identity, guidance (a pointwise weak worsening at world
$i$ is accepted), and the gate with its verdict. A mixed pair that fails the
gate is incomparable at world $i$ in both directions.

\paragraph{Conditions on the verdicts.} For every model $X$:
\begin{enumerate}[label=(V\arabic*),leftmargin=*,itemsep=2pt]
\item (continuity) $p\mapsto\Phi_i(C_p,X)$ is continuous on $[0,1]$;
\item (signs) $\Phi_i(C_p,X)<0$ whenever $T^\on_i(C_p,X)$ is null and
$T^\off_i(C_p,X)$ is non-null; $\Phi_i(C_p,X)\ge0$ whenever
$T^\off_i(C_p,X)$ is null;
\item (covering) for every $p\in(0,1)$, $\Phi_1(C_p,X)\ge0$ or
$\Phi_2(C_p,X)\ge0$.
\end{enumerate}
(V2) concerns only how the verdict extends to pure comparisons; the rule
itself never consults $\Phi_i$ on them.

\begin{lemma}\label{lem:acyclic}
Fix $i$. There is no cycle $X_0\preceq_iX_1\preceq_i\dots\preceq_iX_m=X_0$ with
$X_j\ne X_{j+1}$ for some $j$. In particular $\preceq_i$ satisfies (TC),
finite-chain consistency, and antisymmetry.
\end{lemma}
\begin{proof}
Call a step $X_j\to X_{j+1}$ nontrivial if $X_j\ne X_{j+1}$. For a nontrivial
accepted step, $T^\on_j\coloneqq T^\on_i(X_j,X_{j+1})$ is non-null: under the
guidance clause because $U_i(X_j)\subsetneq U_i(X_{j+1})$ (strictly, as the
models differ and the ranks determine the advice at each world), and under the
gate by definition. Since level sets of $k_i$ are null and $T^\on_j$ is non-null,
\[
 \beta_j\coloneqq\essinf_{T^\on_j}k_i<\esssup_{T^\on_j}k_i\le
 \essinf_{T^\off_j}k_i\eqqcolon\gamma_j ,
\]
with $\gamma_j=+\infty$ if $T^\off_j$ is null. Let $\beta^*\coloneqq\min_j\beta_j$
over nontrivial steps, attained at $j^*$, and $\gamma\coloneqq\min_j\gamma_j$;
then $\gamma>\beta^*$, because $\gamma_j>\beta_j\ge\beta^*$ for every $j$. Put
$W\coloneqq\{\beta^*<k_i<\gamma\}$. No step removes a step of $W$, so
$U_i(X_j)\cap W$ is nondecreasing in $j$. Step $j^*$ adds the set
$T^\on_{j^*}\cap W$, which is non-null because $\essinf_{T^\on_{j^*}}k_i=\beta^*<\gamma$
and the level set $\{k_i=\beta^*\}$ is null. Hence
$U_i(X_0)\cap W\subsetneq U_i(X_m)\cap W$ up to a non-null set, contrary to
$X_m=X_0$.
\end{proof}

\begin{lemma}\label{lem:sep}
Let $p\in[0,1]$ and let $X$ be any model. If $T^\on_i(C_p,X)$ is non-null,
then $(C_p,X)$ is $i$-separated, for $i=1,2$.
\end{lemma}
\begin{proof}
By Lemma~\ref{lem:threshold}, $T^\on_1(C_p,X)\subseteq I\setminus U_1(C_p)
\subseteq\{k_1\le p\}$ and $T^\off_1(C_p,X)\subseteq U_1(C_p)\subseteq\{k_1\ge
p\}$, up to null sets. World~2 is the same with the threshold $1-p$.
\end{proof}

\begin{theorem}\label{thm:main}
Assume (V1)--(V3). On the full four-advice domain the system
$(\preceq_1,\preceq_2)$ with the selection $p\mapsto C_p$ satisfies:
reflexivity; strict (G); (IB); (TC), and indeed the absence of cycles of
Lemma~\ref{lem:acyclic}; (S): every $C_p$ is coherent and undominated,
competitive on $\supp p$, and every section $A_i(X)\coloneqq\{p:C_p\preceq_iX\}$
is closed; (LA); direct (C). The domain satisfies (LF) and (SC). Every model
other than the $C_p$ is strictly dominated at both worlds by some $C_q$.
\end{theorem}
\begin{proof}
\emph{Reflexivity} is the first clause. \emph{(IB).} The sets $T^\on_i$,
$T^\off_i$, the gate and $\Phi_i$ depend only on the advice of the two models
on the set where they differ. \emph{Strict (G).} If $Y$ is a pointwise weak
world-$i$ improvement of $X$, then $T^\off_i(Y,X)=U_i(Y)\setminus U_i(X)$ is
null and $Y\preceq_iX$. If the improvement is strict on a non-null set, then
$T^\on_i(X,Y)$ is null and $T^\off_i(X,Y)$ is non-null, so no clause gives
$X\preceq_iY$. \emph{(TC)} is Lemma~\ref{lem:acyclic}.

\emph{The sections.} Fix $X$ and $p$, and write $T^\on,T^\off$ for the sets of
$(C_p,X)$ at world $i$. We claim that $C_p\preceq_iX$ iff $\Phi_i(C_p,X)\ge0$.
If $X=C_p$, both hold ((V2) with both sets null gives $\Phi_i\ge0$). If
$T^\off$ is null, both hold, by guidance and (V2). If $T^\on$ is null and
$T^\off$ is non-null, both fail, by strict (G) and (V2). Otherwise
$(C_p,X)$ is $i$-separated (Lemma~\ref{lem:sep}) and the gate clause applies,
which is the claim. So $A_i(X)=\{p:\Phi_i(C_p,X)\ge0\}$, which is closed
by (V1).

\emph{Competitiveness.} For $p\in(0,1)$ use (V3) and the claim. At $p=1$,
$U_1(C_1)=U_1(D_1)$ is null, so $T^\off_1$ is null and $C_1\preceq_1X$; at
$p=0$ the same holds at world~2.

\emph{No model is weakly better than a selected one at both worlds.} Suppose
$X\preceq_1C_p$ and $X\preceq_2C_p$. At world~1, $T^\on_1(X,C_p)\subseteq
U_1(C_p)\subseteq\{k_1\ge p\}$ and $T^\off_1(X,C_p)\subseteq I\setminus
U_1(C_p)\subseteq\{k_1\le p\}$. Separation would put a non-null set of keys
$\ge p$ (essentially, $>p$) below a set of keys $\le p$, which is impossible
unless $T^\off_1$ is null. So $X\preceq_1C_p$ forces $U_1(X)\subseteq U_1(C_p)$,
and likewise $U_2(X)\subseteq U_2(C_p)$: $X$ is pointwise weakly better than
$C_p$ at both worlds. On a mixed problem the rank pairs $(r_1,r_2)$ of the four
advices are $(0,2),(1,1),(2,0),(3,3)$, so no advice is weakly better at both
worlds than a different regular advice; on unanimous problems $C_p$ has ranks
$(0,0)$. Hence $X=C_p$. This gives undominatedness of every $C_p$, hence (LA).
Direct (C) holds for every model, because each world is antisymmetric
(Lemma~\ref{lem:acyclic}): a tie at a single world already forces equality.

\emph{Domination of the rest.} For any $X$, $A_1(X)\ni1$ and $A_2(X)\ni0$ are
closed and cover $[0,1]$, so they meet at some $q$: $C_q\preceq_1X$ and
$C_q\preceq_2X$. If $X\ne C_q$, then, since $C_q\preceq_iX$, antisymmetry
(Lemma~\ref{lem:acyclic}) excludes $X\preceq_iC_q$ at either world, so $C_q$
strictly dominates $X$ at both worlds. So the coherent undominated models are
exactly the $C_p$.

\emph{Domain.} The full domain contains every measurable model, hence every
local variant (LF) and every splice (SC).
\end{proof}

\status{PROVED; confirmed by one independent referee
(\texttt{REFEREE\_gate2d.md}), whose two wording repairs are incorporated.
CHECKED on finite cell models with exact arithmetic, for the instance of
Section~\ref{sec:instance}, by two independent implementations. The referee's
checker covers 75 configurations (three families, including one with a precise
band $l=u$, five cell layouts, five seeds, up to $1024$ models). It verifies
full acyclicity of each world, competitiveness at every critical index against
every model, the section identity, strict (G), and that no selected model is
dominated; there are no failures, and a deliberately invalid family is caught.
Without the gate the same verdict gives $730{,}964$ violations of (TC) on 256
models. Finite cells have atoms, so these checks test the combinatorics, not
closedness; in particular the meeting point used for ``domination of the
rest'' can be missing in finite cells. Scripts: \texttt{gate2d/}.}

\begin{remark}\label{rem:roles}
\emph{Division of labour.} (TC) comes from the gate alone, whatever the
verdict: every accepted comparison adds world-$i$ steps before it removes any
in the priority order, and a cycle would have to remove what it adds at the
bottom of its window. (S) comes from the threshold property of
Lemma~\ref{lem:threshold}: every disagreement with a selected model passes the
gate at both worlds, so each world's verdict applies to every such
disagreement. The verdicts are constrained only by continuity, two sign
conditions and covering. The construction never accepts a comparison that fails the gate,
which is how incomparability organizes the coexistence of the non-additive
comparisons.
\end{remark}

\section{A non-additive instance}\label{sec:instance}

Let $\nu\coloneqq\mu\otimes$counting on $I$, a finite measure. For a pair
$(X,Y)$ let $\on_i\coloneqq\nu(T^\on_i(X,Y))$ and
$\off_i\coloneqq\nu(T^\off_i(X,Y))$, and put
\[
 \Phi_i\coloneqq2\,\on_i-\off_i\qquad(i=1,2).
\]
At each world a newly added rank step costs twice what removing one gains:
this is the author's two-measure rule (design memo, \S7, Proposition~1), with
rank steps in place of error sets and $\mu^+=2\mu^-$, the ``sticky''
direction, which on its own violates (TC) (design memo, \S7).

\begin{proposition}\label{prop:instance}
\begin{enumerate}[label=(\alph*),leftmargin=*,itemsep=2pt]
\item $\Phi_1,\Phi_2$ satisfy (V1)--(V3), so Theorem~\ref{thm:main} applies.
Both worlds accept every comparison out of a selected model that the additive
ranking with weight $\nu$ accepts, and more.
\item No integrable $f>0$ a.e.\ on $I$ endorses the world-1 comparisons out of
selected models, that is, satisfies
$C_p\preceq_1X\Rightarrow\int_{U_1(C_p)}f\,d\nu\le\int_{U_1(X)}f\,d\nu$ for all
$p$ and $X$. The same holds at world~2, and for finitely additive weights
positive on non-null sets. More precisely, for a.e.\ mixed problem $g_0$ the
endorsement already fails for records supported in every neighbourhood of
$g_0$, out of the selected model whose boundary passes through $g_0$. For a
\emph{fixed} boundary point this can fail: some positive field endorses all
records localised near one given boundary point.
\end{enumerate}
\end{proposition}
\begin{proof}
(a) (V1): $\mu(C_p\triangle C_q)\to0$ as $q\to p$, and the sets $T^\on_i$,
$T^\off_i$ of $(C_q,X)$ change by at most three times the set where $C_q$ and
$C_p$ differ, so $\on_i,\off_i$ are continuous in $q$.
(V2): if $T^\on_i$ is null and $T^\off_i$ non-null, $\Phi_i=-\off_i<0$; if
$T^\off_i$ is null, $\Phi_i=2\on_i\ge0$.
(V3): on a mixed problem, a change between two regular advices moves the
world-1 rank by $\delta$ and the world-2 rank by $-\delta$, so it contributes
$|\delta|$ steps to $T^\on_1$ and to $T^\off_2$, or to $T^\off_1$ and to
$T^\on_2$. Every other change of advice out of a selected model (into
contradiction, or on a unanimous problem, where $C_p$ is correct) adds steps
at both worlds and removes none, since $C_p$ is never contradictory. Writing
$a$ and $b$ for the masses of world-1 added and removed steps of the first
kind, and $c_1,c_2\ge0$ for the masses of world-1 and world-2 steps of the
second kind,
\[
 \Phi_1=2(a+c_1)-b,\qquad\Phi_2=2(b+c_2)-a,\qquad\Phi_1+\Phi_2=a+b+2c_1+2c_2\ge0 .
\]
The claim after (a) holds because $2x-y\ge x-y$ for $x\ge0$.

(b) Suppose $f$ endorses at world~1. For a mixed problem write $(t,s)$ for
$t$ and the log-radius; in these coordinates $\mu$ has a positive continuous
density $\varrho$. By the Lebesgue differentiation theorem, a.e.\ point
$g_0=(t_0,s_0)$ is a Lebesgue point of $f(\cdot,1)$ with respect to squares
centred at $g_0$, with $f(g_0,1)>0$; fix such a point and
$p\coloneqq l^{-1}(t_0)$; suppose first $l(p)=t_0<u(p)$. For small $\epsilon>0$ let
$Q_\epsilon\coloneqq\{t_0-\epsilon<t<t_0,\ |s-s_0|<\epsilon\}$ and choose
$R_\epsilon\subseteq\{t_0<t<t_0+3\epsilon,\ |s-s_0|<\epsilon\}$ with
$\mu(R_\epsilon)=2\mu(Q_\epsilon)$, possible for small $\epsilon$ by continuity
of $\varrho$. Both sets lie in the square of side $6\epsilon$ centred at $g_0$
and occupy a fraction of it bounded below. Let $X_\epsilon$ be $C_p$ made
silent on $Q_\epsilon$ (where $C_p$ accepts $g$) and accepting $g$ on
$R_\epsilon$ (where $C_p$ is silent). Then
$T^\on_1(C_p,X_\epsilon)=Q_\epsilon\times\{1\}$ and
$T^\off_1(C_p,X_\epsilon)=R_\epsilon\times\{1\}$, so
$\Phi_1=2\mu(Q_\epsilon)-\mu(R_\epsilon)=0$ and $C_p\preceq_1X_\epsilon$
(Lemma~\ref{lem:sep}). Endorsement gives
$\int_{R_\epsilon}f(\cdot,1)\,d\mu\le\int_{Q_\epsilon}f(\cdot,1)\,d\mu$. At a
Lebesgue point the averages of $f(\cdot,1)$ over $Q_\epsilon$ and $R_\epsilon$
tend to $f(g_0,1)$, so the left side is $(2+o(1))f(g_0,1)\mu(Q_\epsilon)$ and
the right side is $(1+o(1))f(g_0,1)\mu(Q_\epsilon)$, a contradiction. At
world~2 use the boundary $t=u(p)$: let $X'_\epsilon$ be $C_p$ made silent on a
set just above $u(p)$ (where $C_p$ accepts $-g$; world-2 rank $0\to1$, an
added step) and accepting $-g$ on a set of twice the mass just below $u(p)$
(silence in $C_p$; world-2 rank $1\to0$, a removed step). Then $\Phi_2=0$,
$C_p\preceq_2X'_\epsilon$, and the same estimate contradicts endorsement.

If $l(p)=u(p)=t_0$, there is no band: $C_p$ accepts $g$ below $t_0$ (world-1
rank $0$) and $-g$ above (rank $2$). Fix a common Lebesgue point of
$f(\cdot,1)$ and $f(\cdot,2)$. Let $X$ be $C_p$ made silent on $Q_\epsilon$
below $t_0$ (adds $(g,1)$ steps, mass $m$) and accepting $g$ on a set
$R_\epsilon$ above $t_0$ with $\mu(R_\epsilon)=\mu(Q_\epsilon)$ (removes $(g,1)$
and $(g,2)$ steps, mass $2m$). Then $\Phi_1=2m-2m=0$, so $C_p\preceq_1X$, and
endorsement requires $\int_{R_\epsilon}(f(\cdot,1)+f(\cdot,2))\,d\mu\le
\int_{Q_\epsilon}f(\cdot,1)\,d\mu$; in the limit $f(g_0,2)\le0$, a
contradiction. World~2 is symmetric.

\emph{Finitely additive weights.} Let $\pi$ be a bounded charge on $I$,
vanishing on null sets and positive on non-null sets, and write
$\pi=\rho\,\nu+\pi_p$ (Yosida--Hewitt), with $\pi_p$ purely finitely additive
and carried, for every $\eta>0$, by a set $Z$ with $\nu(Z)<\eta$. Positivity
forces $\rho>0$ a.e.: if $\rho=0$ on a non-null $A$, take a carrier $Z$ with
$\nu(Z)<\nu(A)$; then $\pi(A\setminus Z)=0$ with $A\setminus Z$ non-null. In
the constructions above, choose a carrier $Z_\epsilon$ of $\pi_p$ whose mass is
a small fraction of the square, and replace $Q_\epsilon$, $R_\epsilon$ by
$Q_\epsilon\setminus Z_\epsilon$ and by a subset of $R_\epsilon\setminus
Z_\epsilon$ with the required mass ratio. These sets still occupy a fraction of
the square bounded below, $\pi_p$ vanishes on them, and the Lebesgue-point
estimate for $\rho$ gives the same contradiction.

\emph{A fixed boundary point.} Fix $p$ with $l(p)<u(p)$, $t_0=l(p)$, and near
$g_0$ put $f(g,1)=1$ for $t>t_0$, $f(g,1)=11$ for $t<t_0$, and
$f(g,2)=f(g,3)=11$. For a record out of $C_p$ supported near $g_0$ and away
from $\{t=u(p)\}$, the removed steps are $(g,1)$ steps above $t_0$, where $f=1$,
and every added step has $f=11$; acceptance gives $\off_1\le2\on_1$, hence
$\int_{T^\off_1}f=\off_1\le11\,\on_1=\int_{T^\on_1}f$. So the non-endorsement is
a property of almost every boundary point, not of each one.
\end{proof}

\noindent Fields that vanish on some non-null set can endorse: for example
$f=1$ on steps with $j=3$ and $f=0$ elsewhere, since selected models have no
such steps. The restriction to $f>0$ a.e.\ is the relevant one.

\status{PROVED; confirmed by one referee with three repairs, all
incorporated (the case $l(p)=u(p)$, finitely additive weights, and ``almost
every'' in place of ``any'' boundary point). Part (b) is the FF2-type
non-endorsement property, now inside a system satisfying every axiom. Since
$l^{-1}$ need not preserve null sets, ``a.e.\ $t_0$'' yields a dense set of
indices $p$, in particular indices arbitrarily close to the target.}

\begin{remark}\label{rem:compare}
\emph{Relation to earlier rules.}
\begin{itemize}[leftmargin=*,itemsep=2pt]
\item Proposition~1 failed because comparisons $X\preceq Y\preceq Z$ with
$Z$ a strict improvement of $X$ could be strung together (design memo, \S7: $X$
with errors on $Q\cup R$, $Y$ on $P\cup R$, $Z$ on $Q$). The first step adds
steps on $P$ and removes them on $Q$; the second adds on $Q$ and removes on
$P\cup R$. Both pass the gate only if the keys of $P$ lie below those of $Q$
and the keys of $Q$ below those of $P$, which is impossible for non-null sets.
\item Neither world is complete, and neither adds up
(Proposition~\ref{prop:instance}(b), including finitely additive weights).
Neither accepted set is a convex cone of records: two gated records whose keys
are interleaved (adds on $Q_1$, removes on $R_1$, adds on $Q_2$, removes on
$R_2$, with $Q_1<R_1<Q_2<R_2$ in key) have a sum that fails the gate.
\item \emph{Leading charges} (PROVED, with the write-up's definition transposed
to rank steps: $\pi(B\setminus A)>\pi(A\setminus B)$ forces $A\le_iB$, and
$A\le_iB$ forces $\pi(A)\le\pi(B)$). Neither world has a nonzero leading charge.
If $\pi\ne0$, pick a non-null $S$ inside one kind of step (fixed $j$, mixed or
unanimous problems) with $\pi(S)>0$, and $c$ with $S\cap\{k<c\}$ and
$S\cap\{k>c\}$ non-null. Split off a non-null $D\subseteq S\cap\{k<c\}$ with
$\pi(D)<\pi(S)/2$, and put $P\coloneqq S\setminus D$. Let $X$ and $Y$ agree
except that $X$ has world-$i$ rank $j-1$ on the problems under $P$ and $j$ on
those under $D$, and $Y$ the reverse. Then $T^\on=P$, $T^\off=D$ and
$\pi(P)>\pi(D)$, so the first property demands $X\preceq_iY$; but $P$ has keys
above those of $D$, the gate fails, and $X\not\preceq_iY$. This rests only on
incompleteness, so it makes the comparison with the 15 October status note
weaker than it looks. That note recorded that no world in a cover was known
which is incomplete, fails to add up, is not CP-a, has no leading charge, and
has certificates that are not point-localised. The worlds here meet the first
four conditions, in two worlds only; the last has not been checked. A non-strict
purely finitely additive $\{0,1\}$-valued charge concentrated at the bottom of
the key order does satisfy the second property: it is the leading criterion of
the lexicographic order in the next item.
\item Each world is contained in the lexicographic order with a
\emph{continuum} of keys, low keys first: $X$ is strictly better than $Y$ when
some $\{k<c\}$ meets $T^\on_i(X,Y)$ in a non-null set and $T^\off_i(X,Y)$ in a
null set. It is not lexicographic over countably many additive criteria, but
only for the trivial reason that such orders are complete and these worlds are
not. Whether the write-up's lexicographic exclusions apply to sub-relations of
a lexicographic order has not been examined.
\item Any verdicts meeting (V1)--(V3) work, including verdicts selected from
the disagreement in the sense of the five-region construction, provided they
are continuous along the family. Exact patterns of the five-region kind can
be added to the accepted comparisons of a world without affecting
Theorem~\ref{thm:main}, provided they (1) pass the gate, (2) respect (IB), that
is, are closed under changes off the difference set, and (3) add, for each
model $X$, a closed set of indices to $A_i(X)$. (TC) is then unaffected by
Lemma~\ref{lem:acyclic}, and undominatedness is unaffected because no gated
comparison into a selected model exists.
\end{itemize}
\end{remark}

\section{Why two worlds are special}\label{sec:three}

The two-world laboratory has been solved too cheaply to discriminate between
verdicts: the gate makes (TC) automatic and nesting makes (S) automatic. The
next proposition locates exactly what fails with three worlds.

Let $n=3$, let the domain be the full four-advice domain, and suppose world $i$
uses a gate of the form in Definition~\ref{def:rule} for some key function
$k_i$ on its rank steps with null level sets: it accepts only identity,
guidance, and $i$-separated pairs. Call $(a,b)$ a \emph{crossed pair for $C$ at
world $i$} if $a\subseteq I\setminus U_i(C)$ and $b\subseteq U_i(C)$ are
non-null and $\esssup_bk_i<\essinf_ak_i$: a step that $C$ lacks has higher key
than a step that $C$ has.

\begin{proposition}\label{prop:three}
Let $p$ be an interior index. Suppose that for each $i\in\{1,2,3\}$ there is a
crossed pair $(a_i,b_i)$ for $C_p$ at world $i$, such that the six sets of
problems underlying $a_1,b_1,a_2,b_2,a_3,b_3$ are pairwise disjoint. Then $C_p$
is not competitive on $\{1,2,3\}$.
\end{proposition}
\begin{proof}
Let $X$ be $C_p$ with contradictory advice on the problems under
$a_1,a_2,a_3$ and with the advice of $D_i$ on the problems under $b_i$. At
world $i$, contradiction has rank $3$, so $T^\on_i(C_p,X)\supseteq a_i$; the
advice of $D_i$ has world-$i$ rank $0$, so $T^\off_i(C_p,X)\supseteq b_i$.
Hence $\esssup_{T^\on_i}k_i\ge\essinf_{a_i}k_i>\esssup_{b_i}k_i\ge
\essinf_{T^\off_i}k_i$: the pair fails the gate at world $i$. It is not a
guidance pair, because $T^\off_i$ is non-null, and $X\ne C_p$. So
$C_p\not\preceq_iX$ for $i=1,2,3$.
\end{proof}

\status{PROVED; confirmed by the referee.}

The disjointness hypothesis is no restriction: non-null subsets of a crossed
pair still cross, and $\mu$ is atomless, so any three crossed pairs, one per
world, can be shrunk onto pairwise disjoint sets of problems. Also, ``no
crossed pair at world $i$'' is equivalent to ``$U_i(C)$ is a threshold set
$\{k_i\ge c\}$ up to null sets''. So under a strong gate every interior selected
model must be a threshold set at some world. The threshold sets of one key
order are nested, so the interior selected models that are world-$i$ threshold
sets form a chain under pointwise world-$i$ dominance. For precise models
$\Dscr_{\{p\}}$ and $\Dscr_{\{q\}}$, pointwise weak world-1 dominance holds
only if $p$ lies on the segment from $q$ to $e_1$ (or conversely), so such a
chain is one-dimensional. A selection through an interior imprecise target
whose image is genuinely two-dimensional cannot be covered by three such
chains. (CONJECTURE: nothing proved here shows that (S) forces a
two-dimensional image. A selection whose interior image lies in three chains,
for instance in a tripod, is not excluded here; the earlier result that the
selected profile cannot factor continuously through a tree bears on the tripod,
but its hypotheses have not been checked against this setting.)

\begin{remark}\label{rem:triangle}
\emph{What a weaker gate must respect in three worlds.} Drop the gate and let
world $i$ judge two-point swaps by a scale-free rule: accept ``add steps on a
small set at $x$ (mass $m_x$), remove steps on a small set at $y$ (mass
$m_y$)'' iff $m_x\ge\rho(x,y)\,m_y$. Take three locations $a,b,c$ and suppose
the selected models' world-$i$ step sets cross cyclically in \emph{both}
orientations: some selected model has $b$ but not $a$, one has $c$ but not
$b$, one has $a$ but not $c$, and likewise with the orientation reversed. This
needs non-nested step sets, so it cannot occur in two worlds; it occurs for
precise families in three. By (IB) the swaps then act on three local models
$W_1,W_2,W_3$ (steps on $\{b,c\}$, $\{c,a\}$, $\{a,b\}$, the same elsewhere),
and every swap between them is judged by $\rho$ because its local advice
matches a selected model. Write $P_f=\rho(a,b)\rho(b,c)\rho(c,a)$ and
$P_r=\rho(a,c)\rho(c,b)\rho(b,a)$. (TC) forces $P_f\ge1$ and $P_r\ge1$.
\emph{Proof} (supplied by the referee). If $P_f<1$, the forward links
$W_1\preceq W_2\preceq W_3\preceq W_1$ hold on a non-empty open set of mass
ratios. To avoid a violation every reverse link must hold on all of it, which
gives $\rho(a,c)\le\rho(a,b)\rho(b,c)$ and its two rotations, hence $P_r\le
P_f^2<1$. Symmetrically $P_f\le P_r^2$, so $P_r\le P_r^4$, which is impossible
for $0<P_r<1$. With $\rho\equiv\frac12$ (the sticky rule) and masses
$(m_a,m_b,m_c)=(1,2,3)$ the violation is explicit: $W_1\preceq W_2\preceq
W_3\preceq W_1$ and $W_1\not\preceq W_3$. \emph{Consequence.} The sticky
direction cannot sit in two-point swaps around such triangles. Non-additive
spreads in the other direction ($\rho\ge1$, for instance $\rho\equiv2$) are
triangle-safe, but every swap whose mass ratio lies in
$(\rho(y,x)^{-1},\rho(x,y))$ is then accepted in neither direction at world
$i$, and must be defended at another world.
\end{remark}

\status{PROVED under the double-crossing hypothesis (argument supplied by the referee; one reader). That precise families in three worlds are doubly crossing is CHECKED by a half-plane construction, not written out.}

\paragraph{Next steps.}
\begin{enumerate}[leftmargin=*,itemsep=2pt]
\item Three worlds. A fixed linear priority order per world fails at every
interior index whose selected model is a threshold set at no world
(Proposition~\ref{prop:three}); that this rules out fixed orders altogether is
CONJECTURE. Sticky verdicts on two-point swaps around
doubly crossing triples are excluded (Remark~\ref{rem:triangle}). Two relaxations remain to be designed and
screened: (a) priority orders chosen from the disagreement itself, with (TC)
checked directly rather than inherited from one lexicographic order; (b)
envelopes with a bid--ask spread on swaps, with the non-additivity in cycles
of length at least four and the spread zones defended at another world. A
cyclic key order with short arcs gives (TC) for local records (three arcs of
less than a third of the circle never cover it), but it rejects every record
spread around the circle, so on its own it cannot supply (S).
\item Screen both against FF2 Thm~4.1, S3 and S9 before any proving round.
\end{enumerate}

\end{document}
